```latex
\chapter{Projeto e Planos de Trabalho - Enzo Ribas}
\label{cap:Projeto_e_Planos_de_Trabalho}

\section{Projeto}

\noindent
Dados do projeto de pesquisa

\begin{description}
  \item[Código:] PIC704-2024

  \item[Título:] Comparação da Implementação de Métodos Numéricos em Diferentes Paradigmas de Programação

  \item[Orientador:] Luis Alberto D'Afonseca
  \item[Co-orientador:] Lucas Martins Rocha 

  \item[Discente:] Enzo Rocha Leite Diniz Ribas

  \item[Tipo de Bolsa:] Voluntário (IC)

  \item[Direcionamento:] Iniciação Científica

  \item[Status:] Semana C\&T Pendente

  \item[Edital:] Edital DPPG nº 73/2024 -- PIBIC FAPEMIG

  \item[Cota:] PIBIC FAPEMIG 2025--2026
               (28/01/2025 a 28/02/2026)
\end{description}

\noindent
Área de Conhecimento

\begin{description}
  \item[Grande Área:] Ciências Exatas e da Terra

  \item[Área:] Matemática

  \item[Subárea:] Matemática Aplicada

  \item[Especialidade:] Análise Numérica
\end{description}

%------------------------------------------------------------------------------%
\section{Plano de Trabalho}

\noindent
\textbf{Título:} Implementação de métodos numéricos em linguagens funcionais

%------------------------------------------------------------------------------%
\subsection{Introdução/Justificativa}

Os métodos numéricos são ferramentas essenciais nas áreas de engenharia e
ciências, onde desempenham um papel fundamental na modelagem de sistemas
complexos. Textos e publicações especializadas destacam a importância dessas
técnicas no contexto computacional, como no livro
\textit{Applied Numerical Methods with MATLAB for Engineers and Scientists},
de Steven Chapra, que enfatiza aplicações práticas. Clássicos como
\textit{Numerical Methods for Scientists and Engineers}, de Richard Hamming
(1962), estabeleceram fundamentos que permanecem relevantes, e
\textit{Numerical Analysis}, de R. L. Burden e J. D. Faires, continua sendo
uma referência atual, abrangendo métodos clássicos e contemporâneos da análise
numérica.

Os paradigmas de programação, por sua vez, oferecem diferentes abordagens para
a construção de software, influenciando diretamente como métodos numéricos são
implementados e executados. Cada paradigma -- seja o imperativo, orientado a
objetos ou funcional -- apresenta suas particularidades em termos de estrutura
de código, mutabilidade de variáveis e organização dos processos
computacionais. A escolha do paradigma afeta aspectos como legibilidade,
desempenho e facilidade de manutenção do código, aspectos cruciais para o
desenvolvimento de soluções numéricas eficientes.

%------------------------------------------------------------------------------%
\subsection{Objetivos}

O objetivo deste projeto é investigar e comparar a implementação de métodos
numéricos em diferentes paradigmas de programação, analisando como cada
abordagem influencia aspectos essenciais do desenvolvimento e execução desses
métodos. Para isso, os métodos selecionados serão implementados em linguagens
representativas dos paradigmas imperativo, orientado a objetos e funcional.

Através dessa análise, o projeto busca identificar as vantagens e desvantagens
de cada paradigma em termos de eficiência, legibilidade, facilidade de
manutenção e expansão do código, fornecendo uma visão mais clara sobre o
impacto das escolhas de paradigma na resolução de problemas numéricos.

Pretendemos que, ao realizar este projeto, os alunos adquiram uma compreensão
prática e teórica dos principais paradigmas de programação e sua aplicabilidade
em problemas científicos, desenvolvendo habilidades de programação em diversas
linguagens. Eles aprenderão sobre a implementação dos métodos numéricos
fundamentais, além de técnicas de análise de desempenho e complexidade de
código.

O projeto permite que explorem os impactos de diferentes paradigmas em termos
de eficiência, legibilidade e manutenção, promovendo uma visão crítica sobre a
escolha de linguagem e paradigma em cenários de desenvolvimento, especialmente
no contexto de aplicações numéricas e científicas.

%------------------------------------------------------------------------------%
\subsection{Método Científico}

A metodologia do projeto envolve a escolha e implementação de métodos
numéricos, como o método de Euler para equações diferenciais ordinárias e o
método de Newton--Raphson para raízes de funções, em diferentes paradigmas de
programação, com o intuito de comparar suas características.

Para isso, serão selecionadas linguagens representativas dos paradigmas
imperativo (C, Octave), orientado a objetos (Java ou Python) e funcional
(Haskell ou F\#). Cada método será implementado em cada linguagem respeitando
as boas práticas e os princípios do respectivo paradigma, assegurando uma
comparação coerente e representativa.

Em seguida, as implementações serão avaliadas quanto ao desempenho
(tempo de execução e uso de memória), complexidade do código
(linhas de código e estrutura algorítmica) e legibilidade e manutenção
(modularidade e clareza). Com esses dados, será feita uma análise comparativa
para identificar as vantagens e desvantagens de cada paradigma na implementação
de métodos numéricos.

%------------------------------------------------------------------------------%
\subsection{Habilidades Adquiridas}

Este projeto visa proporcionar ao estudante uma formação sólida em métodos
numéricos e modelagem matemática, capacitando-o a seguir com pesquisa de nível
de mestrado. Os conhecimentos matemáticos adquiridos serão essenciais para o
desenvolvimento do projeto e contribuirão significativamente para a compreensão
de conteúdos abordados nas disciplinas matemáticas dos cursos de engenharia,
fortalecendo a base acadêmica do aluno.

De forma mais específica, ao realizar este projeto, o aluno desenvolverá uma
compreensão prática e teórica dos principais paradigmas de programação e sua
aplicação em problemas científicos. Ele adquirirá habilidades em diferentes
linguagens de programação e aprenderá a implementar métodos numéricos
fundamentais, além de dominar técnicas de análise de desempenho, complexidade
de código e otimização de soluções.

Esse processo proporcionará ao estudante um entendimento mais profundo sobre
como a escolha do paradigma de programação influencia a eficiência e a
manutenção de soluções numéricas em contextos reais.

Em particular, este plano de trabalho enfatiza as linguagens funcionais. Essas
linguagens são pouco utilizadas, apesar de, em princípio, oferecerem vantagens
sobre as procedurais. Neste projeto, pretendemos explorar essas características
e avaliar se as vantagens são significativas.

%------------------------------------------------------------------------------%
\subsection{Referências}

\begin{enumerate}
  \item CHAPRA, S. C.; CANALE, R. P.
        \textit{Applied Numerical Methods with MATLAB for Engineers and Scientists}.
        McGraw-Hill Education, 2015.

  \item BURDEN, R. L.; FAIRES, J. D.
        \textit{Numerical Analysis}.
        Brooks Cole, 2010.

  \item GABBRIELLI, M.; MARTINI, S.; GIALLorenzo, S.
        \textit{Programming Languages: Principles and Paradigms}.
        Germany: Springer International Publishing, 2023.

  \item MELO, A. C. V. d.; SILVA, F. S. C. d.
        \textit{Princípios de linguagens de programação}.
        Brazil: Editora Blucher, 2003.

  \item STRUIK, D. J.
        \textit{A Source Book in Mathematics, 1200--1800}.
        Princeton University Press, 1987.

  \item HAMMING, R. W.
        \textit{Numerical Methods for Scientists and Engineers}.
        Dover Publications, 1962.

  \item RUGGIERO, M. A. G.; DA ROCHA LOPES, V. L.
        \textit{Cálculo numérico: aspectos teóricos e computacionais}.
        Pearson/Makron, 1998.
\end{enumerate}

%------------------------------------------------------------------------------%
% Cronograma de atividades
%------------------------------------------------------------------------------%

\section*{Corpo do Plano de Trabalho}

\subsection*{Título}

\begin{center}
    \textbf{
    Implementação de métodos numéricos em linguagens funcionais
    }
\end{center}

% ------------------------------------------------------------
% Introdução
% ------------------------------------------------------------
\subsection*{Introdução e Justificativa}

Os métodos numéricos são ferramentas essenciais nas áreas de engenharia e
ciências, onde desempenham um papel fundamental na modelagem de sistemas
complexos. Textos e publicações especializadas destacam a importância dessas
técnicas no contexto computacional, como no livro
\textit{Applied Numerical Methods with MATLAB for Engineers and Scientists},
de Steven Chapra, que enfatiza aplicações práticas. Clássicos como
\textit{Numerical Methods for Scientists and Engineers}, de Richard Hamming
(1962), estabeleceram fundamentos que permanecem relevantes, e
\textit{Numerical Analysis}, de R. L. Burden e J. D. Faires, continua sendo
uma referência atual, abrangendo métodos clássicos e contemporâneos da análise
numérica.

Os paradigmas de programação, por sua vez, oferecem diferentes abordagens
para a construção de software, influenciando diretamente como métodos
numéricos são implementados e executados. Cada paradigma -- seja o
imperativo, orientado a objetos ou funcional -- apresenta suas particularidades
em termos de estrutura de código, mutabilidade de variáveis e organização dos
processos computacionais. A escolha do paradigma afeta aspectos como
legibilidade, desempenho e facilidade de manutenção do código, aspectos
cruciais para o desenvolvimento de soluções numéricas eficientes.

% ------------------------------------------------------------
% Objetivos
% ------------------------------------------------------------
\subsection*{Objetivos}

O objetivo deste projeto é investigar e comparar a implementação de métodos
numéricos em diferentes paradigmas de programação, analisando como cada
abordagem influencia aspectos essenciais do desenvolvimento e execução desses
métodos.

Para isso, os métodos selecionados serão implementados em linguagens
representativas dos paradigmas imperativo, orientado a objetos e funcional.
Através dessa análise, o projeto busca identificar as vantagens e desvantagens
de cada paradigma em termos de eficiência, legibilidade, facilidade de
manutenção e expansão do código, fornecendo uma visão mais clara sobre o
impacto das escolhas de paradigma na resolução de problemas numéricos.

Pretendemos que, ao realizar este projeto, os alunos adquiram uma compreensão
prática e teórica dos principais paradigmas de programação e sua aplicabilidade
em problemas científicos, desenvolvendo habilidades de programação em diversas
linguagens. Eles aprenderão sobre a implementação dos métodos numéricos
fundamentais, além de técnicas de análise de desempenho e complexidade de
código.

O projeto permite que explorem os impactos de diferentes paradigmas em termos
de eficiência, legibilidade e manutenção, promovendo uma visão crítica sobre a
escolha de linguagem e paradigma em cenários de desenvolvimento, especialmente
no contexto de aplicações numéricas e científicas.

% ------------------------------------------------------------
% Metodologia
% ------------------------------------------------------------
\subsection*{Metodologia}

A metodologia do projeto envolve a escolha e implementação de métodos
numéricos, como o método de Euler para EDOs e o método de Newton--Raphson para
raízes de funções, em diferentes paradigmas de programação, com o intuito de
comparar suas características.

Para isso, serão selecionadas linguagens representativas dos paradigmas
imperativo (C, Octave), orientado a objetos (Java ou Python) e funcional
(Haskell ou F\#). Cada método será implementado em cada linguagem respeitando
as boas práticas e os princípios do respectivo paradigma, assegurando uma
comparação coerente e representativa.

Em seguida, as implementações serão avaliadas quanto ao desempenho
(tempo de execução e uso de memória), complexidade do código
(linhas de código e estrutura algorítmica) e legibilidade e manutenção
(modularidade e clareza).

Com esses dados, será feita uma análise comparativa para identificar as
vantagens e desvantagens de cada paradigma na implementação de métodos
numéricos.

% ------------------------------------------------------------
% Habilidades
% ------------------------------------------------------------
\subsection*{Habilidades Adquiridas}

Este projeto visa proporcionar ao estudante uma formação sólida em métodos
numéricos e modelagem matemática, capacitando-o a seguir com pesquisa de nível
de mestrado. Os conhecimentos matemáticos adquiridos serão essenciais para o
desenvolvimento do projeto e contribuirão significativamente para a compreensão
de conteúdos abordados nas disciplinas matemáticas dos cursos de engenharia,
fortalecendo a base acadêmica do aluno.

De forma mais específica, ao realizar este projeto, o aluno desenvolverá uma
compreensão prática e teórica dos principais paradigmas de programação e sua
aplicação em problemas científicos. Ele adquirirá habilidades em diferentes
linguagens de programação e aprenderá a implementar métodos numéricos
fundamentais, além de dominar técnicas de análise de desempenho, complexidade
de código e otimização de soluções.

Esse processo proporcionará ao estudante um entendimento mais profundo sobre
como a escolha do paradigma de programação influencia a eficiência e a
manutenção de soluções numéricas em contextos reais.

Em particular, este plano de trabalho enfatiza as linguagens funcionais. Essas
linguagens são pouco utilizadas, apesar de, em princípio, oferecerem vantagens
sobre as procedurais. Neste projeto, pretendemos explorar essas características
e avaliar se as vantagens são significativas.

% ------------------------------------------------------------
% Referências
% ------------------------------------------------------------
\subsection*{Referências}

\begin{enumerate}[label=\arabic*., leftmargin=*]

    \item CHAPRA, S. C.; CANALE, R. P.
    \textit{Applied Numerical Methods with MATLAB for Engineers and Scientists}.
    McGraw-Hill Education, 2015.

    \item BURDEN, R. L.; FAIRES, J. D.
    \textit{Numerical Analysis}.
    Brooks Cole, 2010.

    \item GABBRIELLI, M.; MARTINI, S.; GIAllorenzo, S.
    \textit{Programming Languages: Principles and Paradigms}.
    Germany: Springer International Publishing, 2023.

    \item MELO, A. C. V. d.; SILVA, F. S. C. d.
    \textit{Princípios de linguagens de programação}.
    Brazil: Editora Blucher, 2003.

    \item STRUIK, D. J.
    \textit{A Source Book in Mathematics, 1200--1800}.
    Princeton University Press, 1987.

    \item HAMMING, R. W.
    \textit{Numerical Methods for Scientists and Engineers}.
    Dover Publications, 1962.

    \item RUGGIERO, M. A. G.; DA ROCHA LOPES, V. L.
    \textit{Cálculo numérico: aspectos teóricos e computacionais}.
    Pearson/Makron, 1998.

\end{enumerate}

% ------------------------------------------------------------
% Cronograma
% ------------------------------------------------------------
\clearpage
```